% Part: many-valued-logic
% Chapter: sequent-calculus
% Section: rules-and-proofs

\documentclass[../../../include/open-logic-section]{subfiles}

\begin{document}

\olfileid{mvl}{seq}{rul}

\olsection{القواعد وال\usetoken{P}{derivation}}

نجعل~${\OLId{greek-variable}{Gamma}}, {\OLId{greek-variable}{Delta}}, \Pi, \Lambda$ فيما يأتي أسماء لمتتاليات
منتهية من ال!!{sentence}s، وندخل في ذلك المتتالية الخالية.

\begin{defn}[التتابعية]
\emph{التتابعية ذات الجوانب~${\OLId{latin-ordinary}{n}}$} تعبير من الصورة
\[
{\OLId{greek-variable}{Gamma}}_{\OLMathNumeral{1}} \nSequent \dots \nSequent {\OLId{greek-variable}{Gamma}}_{\OLId{latin-ordinary}{n}}
\]
حيث كل~${\OLId{greek-variable}{Gamma}}_{\OLId{latin-ordinary}{i}}$ متتالية منتهية (وربما خالية) من !!{sentence}s
اللغة~$\Lang L$.
\end{defn}

\begin{defn}[التتابعية الابتدائية]
\emph{التتابعية الابتدائية ذات الجوانب~${\OLId{latin-ordinary}{n}}$} تتابعية ذات جوانب~${\OLId{latin-ordinary}{n}}$ من
الصورة~$!A \nSequent \dots \nSequent !A$ لأي !!{sentence}~$!A$ في
اللغة.

وقد تشتمل اللغة على رابط صفري الرتبة~$\star$، وهو ثابت من ثوابت
القضايا. وحينئذ نعدّ التتابعية~$\dots \nSequent \star \nSequent \dots$
ابتدائية أيضًا، على أن يوضع~$\star$ وحده في الموضع الموافق لقيمة
الصدق~$\tf{\star} \in {\OLId{latin-looped}{V}}$، وتخلو سائر المواضع.
\end{defn}

إذا كان المنطق ذا~${\OLId{latin-ordinary}{n}}$ من القيم ورمزه~$\Log{L}$، خُصّ كل رابط فيه
بقاعدة منطقية لكل قيمة صدق يجوز أن يأخذها في~$\Log{L}$.
وأما ال!!^{derivation}s في حساب التتابعيات ذي الجوانب~${\OLId{latin-ordinary}{n}}$ للمنطق~$\Log{L}$،
فهي أشجار رؤوسها العليا تتابعيات ابتدائية. وكل تتابعية تقع تحت واحدة
أو أكثر إنما تصح في الشجرة إذا نتجت مما فوقها بتطبيق صحيح لإحدى
قواعد روابط~$\Log{L}$.

\begin{defn}[المبرهنات]
تكون !!^a{sentence}~$!A$ \emph{مبرهنة} في منطق ذي~${\OLId{latin-ordinary}{n}}$ من القيم
$\Log{L}$ إذا وُجد !!a{derivation} للتتابعية ذات الجوانب~${\OLId{latin-ordinary}{n}}$ التي
تحتوي~$!A$ وحدها في كل موضع يقابل قيمة صدق مميّزة في~$\Log{L}$،
وتخلو سائر المواضع. ونكتب
$\Proves[\Log{L}] !A$ إذا كانت~$!A$ مبرهنة، ونكتب
$\Proves/[\Log{L}] !A$ إذا لم تكن كذلك.
\end{defn}

\begin{defn}[!!^{derivability}]
تكون !!^a{sentence}~$!A$ \emph{!!{derivable} من} مجموعة
!!{sentence}s~${\OLId{greek-variable}{Gamma}}$ في منطق ذي~${\OLId{latin-ordinary}{n}}$ من القيم~$\Log{L}$، أي
${\OLId{greek-variable}{Gamma}} \Proves[\Log{L}] !A$، إذا وفقط إذا وُجدت مجموعة جزئية
منتهية~${\OLId{greek-variable}{Gamma}}_{\OLMathNumeral{0}} \subseteq {\OLId{greek-variable}{Gamma}}$ ومتتالية~${\OLId{greek-variable}{Gamma}}_{\OLMathNumeral{0}}'$ من
ال!!{sentence}s في~${\OLId{greek-variable}{Gamma}}_{\OLMathNumeral{0}}$ بحيث يكون للتتابعية الآتية
!!a{derivation}:
\[ \Lambda_{\OLMathNumeral{1}} \nSequent \dots \nSequent \Lambda_{\OLId{latin-ordinary}{n}} \]
حيث تكون~$\Lambda_{\OLId{latin-ordinary}{i}}$ هي~$!A$ إذا كان الموضع~${\OLId{latin-ordinary}{i}}$ يقابل قيمة صدق
مميّزة، وتكون~${\OLId{greek-variable}{Gamma}}_{\OLMathNumeral{0}}'$ إذا كان يقابل قيمة صدق غير مميّزة. وإذا لم
تكن~$!A$ !!{derivable} من~${\OLId{greek-variable}{Gamma}}$، كتبنا
${\OLId{greek-variable}{Gamma}} \Proves/[\Log{L}] !A$.
\end{defn}

ولنعتبر منطق~\L ukasiewicz الثلاثي القيم، فحساب تتابعياته ذو ثلاثة
جوانب. في التتابعية~${\OLId{greek-variable}{Gamma}} \nSequent \Pi \nSequent {\OLId{greek-variable}{Delta}}$
موضع~${\OLId{greek-variable}{Gamma}}$ للقيمة~$\False$، وموضع~${\OLId{greek-variable}{Delta}}$ للقيمة~$\True$،
وموضع~$\Pi$ للقيمة~$\Undef$. وابتدائياته على الصورة
$!A \nSequent !A \nSequent !A$. ثم إن~$\True$ وحدها مميّزة؛
ولهذا كان~${\OLId{greek-variable}{Gamma}} \Proves[\LogLuk[{\OLMathNumeral{3}}]] !A$ مكافئًا لوجود
!!a{derivation} للتتابعية~${\OLId{greek-variable}{Gamma}} \nSequent {\OLId{greek-variable}{Gamma}} \nSequent !A$.
ولو عُدّت~$\Undef$ مميّزة أيضًا، لكان المطلوب !!a{derivation}
للتتابعية~${\OLId{greek-variable}{Gamma}} \nSequent !A \nSequent !A$.

\end{document}
